\chapter{The Gauss Map and Shape Operator}
\label{ch:vii16-gauss-map-shape-operator}

For an oriented regular surface \(S\subset\mathbb R^3\), the unit normal is not merely an auxiliary vector.  Its variation records how the tangent plane turns.  The two central objects are
\[
\boxed{N:S\to S^2,\qquad S_p=-\,dN_p:T_pS\to T_pS.}
\]
The first is the Gauss map and the second is the shape operator (Weingarten map).

\section{Local unit normals}

\begin{definition}[Local unit normal]\label{def:vii16-local-normal}
For a regular parametrization \(X:U\subset\mathbb R^2\to S\), the associated unit normal is
\[
N=\frac{X_u\times X_v}{\|X_u\times X_v\|}.
\]
Changing the order of the coordinates reverses the sign.
\end{definition}

\begin{proposition}[Orthogonality]\label{prop:vii16-normal-orthogonal}
At every point of a regular patch,
\[
\langle N,X_u\rangle=\langle N,X_v\rangle=0,\qquad \langle N,N\rangle=1.
\]
Consequently \(dN_p(v)\perp N(p)\), so \(dN_p(v)\in T_pS=T_{N(p)}S^2\).
\end{proposition}

\begin{proof}
The first two identities follow from the cross product. Differentiating
\(\langle N,N\rangle=1\) in a tangent direction \(v\) gives
\(2\langle dN_p(v),N(p)\rangle=0\).
\end{proof}

\section{Orientability and the global Gauss map}

\begin{definition}[Orientable surface]\label{def:vii16-orientable}
A regular surface is orientable if it admits a smooth global unit normal field \(N:S\to S^2\).
\end{definition}

\begin{proposition}[Two choices on a connected surface]\label{prop:vii16-two-orientations}
If \(S\) is connected and orientable, any two global unit normals inducing the same orientation agree, while the opposite orientation is represented by \(-N\).
\end{proposition}

\begin{warning}[The Gauss map requires an orientation]\label{warn:vii16-gauss-orientation}
A nonorientable surface has no globally defined Gauss map into \(S^2\), although every coordinate patch has a local Gauss map.
\end{warning}

\section{Differential of the Gauss map}

\begin{definition}[Shape operator]\label{def:vii16-shape}
For an oriented surface and \(p\in S\), define
\[
S_p=-dN_p:T_pS\to T_pS.
\]
\end{definition}

\begin{theorem}[Weingarten self-adjointness]\label{thm:vii16-self-adjoint}
The shape operator is self-adjoint with respect to the first fundamental form:
\[
\langle S_pu,v\rangle=\langle u,S_pv\rangle.
\]
\end{theorem}

\begin{proof}
In a parametrization, differentiate
\(\langle N,X_u\rangle=\langle N,X_v\rangle=0\). This gives
\[
\langle N_u,X_v\rangle=\langle N_v,X_u\rangle
\]
because mixed partials \(X_{uv}=X_{vu}\). Since \(S=-dN\), the matrix is symmetric with respect to the metric.
\end{proof}

\section{Coordinate matrix of the shape operator}

Let
\[
I=\begin{pmatrix}E&F\\F&G\end{pmatrix},
\qquad
II=\begin{pmatrix}e&f\\f&g\end{pmatrix},
\]
where
\[
E=\langle X_u,X_u\rangle,\quad F=\langle X_u,X_v\rangle,\quad
G=\langle X_v,X_v\rangle,
\]
and
\[
e=\langle X_{uu},N\rangle,\quad
f=\langle X_{uv},N\rangle,\quad
g=\langle X_{vv},N\rangle.
\]

\begin{theorem}[Matrix relation]\label{thm:vii16-matrix}
In the basis \((X_u,X_v)\),
\[
[S]=I^{-1}II.
\]
Equivalently, for tangent vectors \(u,v\),
\[
II(u,v)=I(Su,v).
\]
\end{theorem}

\begin{proof}
From \(\langle N,X_i\rangle=0\), differentiation gives
\[
-\langle N_i,X_j\rangle=\langle N,X_{ij}\rangle.
\]
Thus the bilinear form associated with \(-dN\) through the metric is exactly \(II\).
\end{proof}

\section{Normal curvature}

\begin{definition}[Normal curvature]\label{def:vii16-normal-curvature}
For a nonzero tangent vector \(v\in T_pS\),
\[
k_n(v)=\frac{II(v,v)}{I(v,v)}
      =\frac{\langle S_pv,v\rangle}{\langle v,v\rangle}.
\]
\end{definition}

\begin{theorem}[Meusnier normal-curvature principle]\label{thm:vii16-meusnier}
If a regular curve on the surface has unit tangent \(T\) at \(p\) and curvature vector \(\kappa n_c\), then the surface-normal component is
\[
k_n(T)=\langle \kappa n_c,N\rangle.
\]
Thus normal curvature depends only on the tangent direction, not on the chosen surface curve realizing it.
\end{theorem}

\section{Principal directions preview}

\begin{theorem}[Spectral theorem for the shape operator]\label{thm:vii16-spectral}
At every point, the self-adjoint map \(S_p\) admits an orthonormal eigenbasis of \(T_pS\). Its eigenvalues are the principal curvatures and its eigenvectors are the principal directions.
\end{theorem}

\begin{remark}
VII/17 develops the invariants built from these eigenvalues. Here the emphasis is the operator itself and its geometric origin as the derivative of the Gauss map.
\end{remark}

\section{Examples}

\begin{example}[Plane]\label{ex:vii16-plane}
For a plane, \(N\) is constant. Hence \(dN=0\) and \(S=0\).
\end{example}

\begin{example}[Sphere]\label{ex:vii16-sphere}
For the outward orientation of the sphere \(S_R^2\),
\[
N(p)=\frac{p}{R},\qquad dN_p(v)=\frac{v}{R},
\qquad S_p=-\frac1R\,\operatorname{id}.
\]
With the inward orientation the signs reverse.
\end{example}

\begin{example}[Cylinder]\label{ex:vii16-cylinder}
For \(X(u,v)=(R\cos u,R\sin u,v)\), the normal is
\(N=(\cos u,\sin u,0)\). The axial direction lies in \(\ker S\), while the circular direction has nonzero normal bending.
\end{example}

\section{Geometric meaning}

The shape operator measures first-order rotation of tangent planes. A direction \(v\) lies in \(\ker S_p\) precisely when the surface normal is stationary to first order along \(v\). The determinant and trace of \(S_p\) will become the Gaussian and mean curvature in VII/17.

\section{Exercises}
\begin{exercise}\label{exr:vii16-01}
Show that \(dN_p(v)\) is tangent to \(S^2\) at \(N(p)\).
\end{exercise}
\begin{hint}
Differentiate \(\langle N,N
angle=1\).
\end{hint}
\begin{solution}
The derivative gives \(2\langle dN_p(v),N(p)
angle=0\), which is exactly tangency to the unit sphere.
\end{solution}

\begin{exercise}\label{exr:vii16-02}
Verify that \(N=(X_u\times X_v)/\|X_u\times X_v\|\) is orthogonal to the patch tangent plane.
\end{exercise}
\begin{hint}
Use the defining property of the cross product.
\end{hint}
\begin{solution}
The cross product is orthogonal to both \(X_u\) and \(X_v\), which span the tangent plane.
\end{solution}

\begin{exercise}\label{exr:vii16-03}
Explain how reversing patch orientation changes the local normal.
\end{exercise}
\begin{hint}
Interchange \(u\) and \(v\).
\end{hint}
\begin{solution}
Since \(X_v\times X_u=-X_u\times X_v\), the unit normal changes from \(N\) to \(-N\).
\end{solution}

\begin{exercise}\label{exr:vii16-04}
Show that if \(N\) is a global unit normal then \(-N\) is another.
\end{exercise}
\begin{hint}
Check smoothness and norm.
\end{hint}
\begin{solution}
Negation preserves smoothness and unit length and reverses the chosen orientation.
\end{solution}

\begin{exercise}\label{exr:vii16-05}
Compute the shape operator of a plane.
\end{exercise}
\begin{hint}
Its normal is constant.
\end{hint}
\begin{solution}
Because \(dN=0\), one has \(S=-dN=0\).
\end{solution}

\begin{exercise}\label{exr:vii16-06}
Compute \(S\) for the outward oriented sphere of radius \(R\).
\end{exercise}
\begin{hint}
Use \(N(p)=p/R\).
\end{hint}
\begin{solution}
Then \(dN_p(v)=v/R\), hence \(S_p(v)=-v/R\).
\end{solution}

\begin{exercise}\label{exr:vii16-07}
How does the sphere shape operator change under reversal of orientation?
\end{exercise}
\begin{hint}
Replace \(N\) by \(-N\).
\end{hint}
\begin{solution}
Then \(S=-d(-N)=dN\), so all eigenvalues change sign.
\end{solution}

\begin{exercise}\label{exr:vii16-08}
For the standard cylinder, identify a direction in \(\ker S\).
\end{exercise}
\begin{hint}
Move parallel to the axis.
\end{hint}
\begin{solution}
The normal is independent of the axial coordinate, so its derivative in the axial direction vanishes.
\end{solution}

\begin{exercise}\label{exr:vii16-09}
Prove the identity \(II(u,v)=I(Su,v)\).
\end{exercise}
\begin{hint}
Differentiate \(\langle N,X_i
angle=0\).
\end{hint}
\begin{solution}
The differentiated orthogonality equations yield \(-\langle dN(u),v
angle=\langle N,D_uv
angle\), which is the second fundamental form.
\end{solution}

\begin{exercise}\label{exr:vii16-10}
Derive \([S]=I^{-1}II\).
\end{exercise}
\begin{hint}
Write the previous bilinear identity in matrix form.
\end{hint}
\begin{solution}
The identity is \(II=I[S]\); multiplying by \(I^{-1}\) gives \([S]=I^{-1}II\).
\end{solution}

\begin{exercise}\label{exr:vii16-11}
Why is \(I\) invertible on a regular patch?
\end{exercise}
\begin{hint}
Use linear independence of \(X_u,X_v\).
\end{hint}
\begin{solution}
Its Gram matrix is positive definite because the tangent vectors are independent.
\end{solution}

\begin{exercise}\label{exr:vii16-12}
Prove self-adjointness of \(S\).
\end{exercise}
\begin{hint}
Use symmetry of \(II\).
\end{hint}
\begin{solution}
Since \(II(u,v)=II(v,u)\) and \(II(u,v)=I(Su,v)\), one gets \(I(Su,v)=I(u,Sv)\).
\end{solution}

\begin{exercise}\label{exr:vii16-13}
Show that eigenvectors belonging to distinct eigenvalues are orthogonal.
\end{exercise}
\begin{hint}
Apply self-adjointness.
\end{hint}
\begin{solution}
If \(Su=\lambda u\), \(Sv=\mu v\), then \((\lambda-\mu)\langle u,v
angle=0\).
\end{solution}

\begin{exercise}\label{exr:vii16-14}
Express normal curvature as a Rayleigh quotient.
\end{exercise}
\begin{hint}
Use \(II(v,v)=I(Sv,v)\).
\end{hint}
\begin{solution}
Thus \(k_n(v)=\langle Sv,v
angle/\langle v,v
angle\).
\end{solution}

\begin{exercise}\label{exr:vii16-15}
Show normal curvature is unchanged when \(v\) is rescaled.
\end{exercise}
\begin{hint}
Both numerator and denominator are quadratic.
\end{hint}
\begin{solution}
Replacing \(v\) by \(cv\) multiplies numerator and denominator by \(c^2\).
\end{solution}

\begin{exercise}\label{exr:vii16-16}
What happens to normal curvature when the surface orientation is reversed?
\end{exercise}
\begin{hint}
The shape operator changes sign.
\end{hint}
\begin{solution}
Therefore every normal curvature changes sign.
\end{solution}

\begin{exercise}\label{exr:vii16-17}
Show that \(S_p=0\) implies the Gauss map has zero differential at \(p\).
\end{exercise}
\begin{hint}
Use the definition.
\end{hint}
\begin{solution}
Since \(S_p=-dN_p\), the two statements are equivalent.
\end{solution}

\begin{exercise}\label{exr:vii16-18}
Interpret \(\ker S_p\) geometrically.
\end{exercise}
\begin{hint}
Think about \(dN_p(v)\).
\end{hint}
\begin{solution}
It consists of directions in which the unit normal is stationary to first order.
\end{solution}

\begin{exercise}\label{exr:vii16-19}
For a unit eigenvector \(e\) of \(S\) with eigenvalue \(k\), compute \(k_n(e)\).
\end{exercise}
\begin{hint}
Insert \(Se=ke\).
\end{hint}
\begin{solution}
The quotient becomes \(\langle ke,e
angle/\langle e,e
angle=k\).
\end{solution}

\begin{exercise}\label{exr:vii16-20}
Why does the spectral theorem apply to \(S_p\)?
\end{exercise}
\begin{hint}
The tangent plane is a Euclidean inner-product space.
\end{hint}
\begin{solution}
It is finite-dimensional and \(S_p\) is self-adjoint.
\end{solution}

\begin{exercise}\label{exr:vii16-21}
Relate the second fundamental coefficients to derivatives of the normal.
\end{exercise}
\begin{hint}
Differentiate normal orthogonality.
\end{hint}
\begin{solution}
One gets \(e=-\langle N_u,X_u
angle\), \(f=-\langle N_u,X_v
angle=-\langle N_v,X_u
angle\), \(g=-\langle N_v,X_v
angle\).
\end{solution}

\begin{exercise}\label{exr:vii16-22}
Show that \(f\) has the two equivalent expressions above.
\end{exercise}
\begin{hint}
Differentiate \(\langle N,X_u
angle\) in \(v\) and \(\langle N,X_v
angle\) in \(u\).
\end{hint}
\begin{solution}
Equality follows from \(X_{uv}=X_{vu}\).
\end{solution}

\begin{exercise}\label{exr:vii16-23}
Explain why the Gauss map is constant on a connected planar patch.
\end{exercise}
\begin{hint}
All tangent planes have the same normal line.
\end{hint}
\begin{solution}
A continuous choice of unit normal cannot jump between antipodal values on a connected patch, hence it is constant.
\end{solution}

\begin{exercise}\label{exr:vii16-24}
State the bridge from \(S_p\) to VII/17.
\end{exercise}
\begin{hint}
Use its two eigenvalues.
\end{hint}
\begin{solution}
VII/17 defines principal curvatures as the eigenvalues, Gaussian curvature as their product, and mean curvature as half their sum.
\end{solution}

\section{Solved problem dossiers}
\begin{problem}[Legacy Gauss-map problem]\label{prob:vii16-01}
Starting from a regular parametrization, construct the local Gauss map and prove its differential takes tangent vectors to the tangent plane.
\end{problem}
\begin{solution}
Normalize \(X_u\times X_v\). Differentiating \(\langle N,N
angle=1\) shows \(dN(v)\perp N\), hence \(dN(v)\in T_pS\).
\end{solution}

\begin{problem}[Legacy shape-operator problem]\label{prob:vii16-02}
Derive the shape operator from the Gauss map and explain the sign convention.
\end{problem}
\begin{solution}
Define \(S=-dN\). The minus sign makes the usual outward sphere carry negative eigenvalues under this convention; reversing orientation reverses \(S\).
\end{solution}

\begin{problem}[Legacy Weingarten problem]\label{prob:vii16-03}
Prove the shape operator is self-adjoint.
\end{problem}
\begin{solution}
Differentiate the two equations \(\langle N,X_u
angle=\langle N,X_v
angle=0\), use equality of mixed partials, and translate the symmetric coefficient matrix into the metric inner product.
\end{solution}

\begin{problem}[Coordinate-matrix problem]\label{prob:vii16-04}
Prove \([S]=I^{-1}II\).
\end{problem}
\begin{solution}
From \(II(u,v)=I(Su,v)\), the matrices satisfy \(II=I[S]\). Since \(I\) is positive definite, it is invertible.
\end{solution}

\begin{problem}[Sphere problem]\label{prob:vii16-05}
Compute \(N,dN,S\) for the sphere of radius \(R\).
\end{problem}
\begin{solution}
For the outward normal, \(N(p)=p/R\); thus \(dN(v)=v/R\) and \(S(v)=-v/R\).
\end{solution}

\begin{problem}[Cylinder problem]\label{prob:vii16-06}
Compute the qualitative eigendirections of the circular cylinder.
\end{problem}
\begin{solution}
The normal varies around circles but not along generators. Hence the axial direction has eigenvalue \(0\) and the circular direction has the nonzero eigenvalue.
\end{solution}

\begin{problem}[Normal-curvature problem]\label{prob:vii16-07}
Prove that normal curvature depends only on a tangent direction.
\end{problem}
\begin{solution}
The surface-normal component of the curve acceleration is \(II(T,T)\), so for a unit tangent it equals \(I(ST,T)\), determined by the tangent vector and the surface.
\end{solution}

\begin{problem}[Orientation problem]\label{prob:vii16-08}
Track \(N,S,II\), and normal curvature under orientation reversal.
\end{problem}
\begin{solution}
Replacing \(N\) by \(-N\) multiplies \(dN,S,II\), and every normal curvature by \(-1\), while \(I\) is unchanged.
\end{solution}

\begin{problem}[Kernel problem]\label{prob:vii16-09}
Characterize directions along which the normal is stationary to first order.
\end{problem}
\begin{solution}
They are exactly \(\ker dN_p=\ker S_p\).
\end{solution}

\begin{problem}[Principal-direction preview]\label{prob:vii16-10}
Use the spectral theorem to produce an orthonormal basis adapted to the shape operator.
\end{problem}
\begin{solution}
Self-adjointness gives an orthonormal eigenbasis of \(T_pS\).
\end{solution}

\begin{problem}[Second-form reconstruction problem]\label{prob:vii16-11}
Show \(I\) and \(S\) determine \(II\).
\end{problem}
\begin{solution}
The formula is \(II(u,v)=I(Su,v)\).
\end{solution}

\begin{problem}[Gauss-map rank problem]\label{prob:vii16-12}
Explain what rank \(0,1,2\) of \(dN_p\) means infinitesimally.
\end{problem}
\begin{solution}
Rank records the number of independent tangent directions in which the normal changes; because \(S=-dN\), it is the rank of the shape operator.
\end{solution}

\section{Challenge problems}
\begin{problem}[Weingarten equations]\label{prob:vii16-ch01}
Derive explicit formulas for \(N_u\) and \(N_v\) in the basis \(X_u,X_v\).
\end{problem}
\begin{solution}
Write \(N_u=aX_u+bX_v\) and pair with \(X_u,X_v\). The resulting linear system is \(I(a,b)^T=-(e,f)^T\). Similarly \(I(c,d)^T=-(f,g)^T\) for \(N_v=cX_u+dX_v\).
\end{solution}

\begin{problem}[Gauss-map local behavior]\label{prob:vii16-ch02}
Suppose \(S_p\) is invertible. What does the inverse function theorem say about the Gauss map near \(p\)?
\end{problem}
\begin{solution}
Since \(dN_p=-S_p\) is an isomorphism between two-dimensional tangent spaces, \(N\) is a local diffeomorphism from the surface to an open subset of \(S^2\).
\end{solution}

\begin{problem}[Umbilic operator criterion]\label{prob:vii16-ch03}
Show that if every tangent direction has the same normal curvature at \(p\), then \(S_p=\lambda I\).
\end{problem}
\begin{solution}
For a self-adjoint operator, equality of the Rayleigh quotient on every unit vector forces both eigenvalues to equal the same number \(\lambda\); hence \(S_p=\lambda I\).
\end{solution}

\begin{problem}[Rank-one shape operator]\label{prob:vii16-ch04}
Assume \(S_p\) has rank one. Describe its kernel and image geometrically.
\end{problem}
\begin{solution}
Self-adjointness gives \(T_pS=\ker S_p\oplus(\ker S_p)^\perp\). The image is the one-dimensional orthogonal complement of the kernel and is the unique direction in which the normal changes to first order.
\end{solution}

\begin{problem}[Coordinate invariance]\label{prob:vii16-ch05}
Explain why \(I^{-1}II\) transforms as an endomorphism under a change of surface coordinates.
\end{problem}
\begin{solution}
Both forms undergo congruence transformations by the Jacobian. Their combination \(I^{-1}II\) changes by similarity, exactly the transformation law of a linear operator, so the shape operator is intrinsic to the embedding and orientation rather than to the chosen coordinates.
\end{solution}

\section*{Bridge to VII/17}
VII/17 turns the spectral data of the shape operator into scalar curvature invariants: principal curvatures, Gaussian curvature, and mean curvature.
